Nguồn: \href{https://www.bbc.com/vietnamese/vietnam-44098055}{BBC}

Một cựu quan chức Ban Dân vận nói với BBC rằng bi kịch "không được lắng nghe" của giáo sư Phan Đình Diệu cũng là của giới trí thức Việt Nam.

Giáo sư Phan Đình Diệu, nhà toán học, khoa học máy tính kỳ cựu của Việt Nam qua đời tại Hà Nội sau một thời gian lâm bệnh, hưởng thọ 83 tuổi, gia đình ông xác nhận.

Giáo sư Diệu là người sáng lập và Chủ tịch đầu tiên của Hội Tin học Việt Nam.

Đường nào cho cải cách giáo dục?

"Ông Diệu được ghi nhận là một trong những người có công lao đầu tiên xây dựng và phát triển ngành tin học tại Việt Nam. Trong cuộc đời mình, ông cũng nổi tiếng với sự chính trực và có những đóng góp tâm huyết về chính sách phát triển khoa học giáo dục nước nhà," theo VietnamNet hôm 13/5.

Thập niên 1990, ông là cựu đại biểu Quốc hội khóa VI nhưng sau đó có tin ông "bị gạt bỏ khỏi danh sách đại biểu Quốc hội" vì tham gia hoạt động phong trào dân chủ, đòi đổi mới chính trị để phát triển đất nước.

'Tiếng nói bị bỏ ngoài tai'
Hôm 14/5, trả lời BBC từ Hà Nội, Giáo sư, Nhà nghiên cứu Nguyễn Khắc Mai, cựu Vụ trưởng Nghiên cứu, Ban Dân vận Trung ương Đảng Cộng sản Việt Nam, nói: "Qua những dịp trò chuyện cùng Giáo sư Phan Đình Diệu, tôi nhận thấy ông là kẻ sĩ chân chính, cương trực."

"Ông có nhiều tư tưởng mới và phát biểu thẳng thắn về những vấn đề dân chủ, không tán thành chế độ độc đảng."

"Đáng tiếc là giới lãnh đạo đã bỏ ngoài tai, không nghe, không sửa."

"Lẽ ra người ta phải đem ý kiến của ông ra phân tích, tranh luận để cho đất nước tốt đẹp hơn."

"Bi kịch không được lắng nghe của Giáo sư Phan Đình Diệu cũng là của giới trí thức Việt Nam."

"Dù ông mất đi, nhưng hy vọng tinh thần Phan Đình Diệu vẫn vĩnh hằng trong giới trí thức chân chính."


Giáo sư Nguyễn Đăng Hưng viết trên trang cá nhân: "Giáo sư Phan Đình Diệu hẳn để lại niềm thương tiếc sâu sắc cho các đồng nghiệp, các học trò cũ, những con người tử tế trong giai đoạn vàng thau lẫn lộn, thế sự không bình thường của đất nước."

"Khi còn tại chức tại Bỉ, trước năm 1990, theo dõi tình trong nước, tôi chú ý đến ông vì những phát biểu khách quan, vô tư của ông về tình hình chính trị, kinh tế, khoa học, xă hội Việt Nam lúc bấy giờ! Trong bối cảnh đơn điệu của những bế tắc triền miên, đây là thái độ dũng cảm hiếm hoi đầy khí phách của một trí thức xă hội chủ nghĩa!"

"Nét thông tuệ đáng nể nhất của ông Diệu là những nhận xét, những quan điểm của ông về tình hình xă hội chính trị Việt Nam và quốc tế. Ông không hề ngần ngại tỏ rõ chính kiến của mình khi có dịp phát biểu thông qua các phương tiện thông tin đại chúng."

'Trong những năm sau 1975, được hỏi về ảnh hưởng chính trị đối với đất nước của một vị quan to nhất lúc bấy giờ, Giáo sư Diệu đã không ngần ngại phát biểu: "Vị lãnh đạo ấy rất vĩ đại, đưa cuộc kháng chiến chống Mỹ tới thắng lợi, nhưng vĩ đại hơn, nếu ông từ chức ngay bây giờ".

"Giáo sư Diệu đã dự đoán không đúng. Vị quan đứng đầu nước ấy làm sao từ chức được và ảnh hưởng về những quyết sách của vị lãnh đạo này trong những năm bao cấp 1975-1986 như thế nào thì nay ai cũng biết."

"Giáo sư Diệu là nhà khoa học, ông không phải là nhà chính trị nhưng những chính kiến công khai và minh bạch của ông về tình hình chính trị xã hội thật là sắc sảo, những phản biện thẳng thắn của ông về chính sách của chính phủ hiện hành thật là thấu đáo thông tuệ," ông Hưng viết.

Về ý thức hệ, trí thức và giáo dục
Hồi năm 1992, khi trả lời Stein Tønnesson, người Na Uy, GS Phan Đình Diệu đã nói về trí thức ở Việt Nam:

"Trái ngược với nền giáo dục thời Pháp, chế độ xã hội chủ nghĩa đào tạo ra những chuyên viên hơn là những trí thức. Chúng tôi có nhiều nhà toán học, nhà vật lý học, nhà sinh học, kỹ sư... và bây giờ thêm nhiều nhà kinh tế. Nhưng chưa bao giờ họ được học để suy nghĩ về các vấn đề của xã hội. Đảng nghĩ hộ cho mọi người. Ý thức chính trị của các chuyên viên nói chung là yếu. Những người giỏi tham gia chính quyền, và tự nhiên là đảng viên.

Rất có thể nhiều chuyên viên, trong cuộc sống riêng, cũng có tư tưởng dân chủ, nhưng không có cách gì kiểm nghiệm điều đó cả.

Thành phần thứ ba là những trí thức được đào tạo trước đây ở miền Nam. Phần đông đã bỏ đi. Tất nhiên có thể họ sẽ trở về giúp nước bằng cách này hay cách khác, nhưng muốn đóng một vai trò chính trị có ý nghĩa, thì người trí thức phải gần gụi nhân dân. Cuộc sống kéo dài ở hải ngoại không phải là mảnh đất thuận lợi cho một lực lượng trí thức tích cực.

Cuối cùng là thanh niên, những người vừa được hay còn đang được đào tạo trong những năm gọi là đổi mới. Những năm gần đây, quả đã có một nền văn nghệ độc lập khởi sắc. Nhưng còn phải có thời gian thì các xu hướng nói trên mới có thể hình thành một lực lượng xã hội chính trị thực sự. Nói tóm lại, kết luận của tôi là hiện nay Việt Nam chưa có một giai cấp trí thức."

Cũng từ khi đó, trong bài trả lời phỏng vấn sau được đăng trên Nordic Newsletter of Asian Studies (1993), ông nói với nhà sử học Na Uy về hệ thống xã hội chủ nghĩa và ý thức hệ Marxist-Leninist như sau:

GS Phan Đình Diệu từng được học giả nước ngoài ví với Andrei Sakharov nổi tiếng của Liên Xô cũ, người sau phản đối nhiều chính sách của Kremlin

" Là một người làm khoa học, đã từ khá lâu tôi phát hiện ra rằng chủ nghĩa Mác-Lê khó có thể mang lại điều gì hữu ích cho một nước muốn thoát ra khỏi nghèo nàn. Nghiên cứu khoa học điện toán, thuyết hệ thống và các vấn đề quản lý hiện đại, tôi nhận ra rằng mô hình xã hội chủ nghĩa như học thuyết Mác-Lê xác định không phù hợp với một nước muốn phát triển về xã hội, kinh tế và khoa học.

Vì những lý do hiển nhiên, Marx và Lenin không có điều kiện tìm hiểu xã hội hiện đại. Song đảng cộng sản ở Việt Nam cũng như ở các nước khác đã ngây thơ tìm cách áp dụng học thuyết của họ vào xã hội hiện đại. Sự thất bại của họ đã được minh chứng rõ rệt nhất trong cuộc cách mạng điện toán, đặc trưng của thế giới trong thập niên 1980.

Sự thất bại của hệ thống kinh tế xã hội chủ nghĩa ngày nay hầu như mọi người đều thấy rõ. Vấn đề còn lại là biết rút ra kết luận một cách đầy đủ và từ bỏ chủ nghĩa Mác-Lê."

Được biết, Stein Tønnesson đã hỏi Phan Đình Diệu khi gặp nhau tại Hà Nội rằng ông có phải là 'Sakharov' của Việt Nam, và được nghe câu trả lời:

"Tôi không phải là một chính khách và tôi không có tham vọng (chính trị). Song, với tư cách một nhà khoa học và một công dân yêu nước, tôi muốn tham gia việc nước. Dân chủ là tham gia việc nước."

Sang năm 2008, trả lời BBC Tiếng Việt, GS Phan Đình Diệu nói về giáo dục Việt Nam:

"Nếu cứ bám chặt quan điểm về giáo dục là phải bồi dưỡng một ý thức hệ nào đó thì sẽ rất khó khăn,"

"Mục tiêu của giáo dục không phải là dạy cho học trò biết đủ mọi thứ".

Theo ông các chính sách về giáo dục ở Việt Nam [cho đến gần 2008] vẫn mang tính chạy theo thành tích bề nổi mà không đề cập đến những vấn đề cơ bản một cách thỏa đáng, hay "chưa trúng đích".

